\section*{Capítulo \ref{ch:g5:air-pressure-weather} --- La presión del aire y el tiempo atmosférico}

\begin{solution}{exo:g5:air-pressure-weather:1}
Se pesa el mismo balón hinchado y después deshinchado: lo único que lo
ha abandonado es \omterm{def:g2:air-around-us:air}{aire}, y la \omterm{def:g3:levers-and-scales:scale}{balanza} marca menos --- la diferencia es
el \omterm{def:g4:weight-and-mass:weight}{peso} del \omterm{def:g2:air-around-us:air}{aire} escapado. El \omterm{def:g2:air-around-us:air}{aire} es materia con \omterm{def:g3:measuring-mass:mass}{masa}.
\end{solution}

\begin{solution}{exo:g5:air-pressure-weather:2}
El empuje del \omterm{def:g2:air-around-us:air}{aire} apretado del fondo de la \omterm{def:g5:air-pressure-weather:pressure}{atmósfera} sobre todas
las superficies que toca. Viene del \omterm{def:g4:weight-and-mass:weight}{peso} del profundo océano de
\omterm{def:g2:air-around-us:air}{aire} que tenemos encima, y empuja desde todas las direcciones ---
hacia abajo, de lado y hacia arriba.
\end{solution}

\begin{solution}{exo:g5:air-pressure-weather:3}
Porque los empujes llegan equilibrados: el \omterm{def:g2:air-around-us:air}{aire} del otro lado de tu
palma --- y el de dentro de tu cuerpo --- empuja de vuelta exactamente
con la misma \omterm{def:g2:pushes-and-pulls:force}{fuerza}. Solo un empuje \emph{desequilibrado}, como el de
la cartulina o el de la botella que se arruga, enseña su poder.
\end{solution}

\begin{solution}{exo:g5:air-pressure-weather:4}
El \omterm{def:g4:weight-and-mass:weight}{peso} del agua empuja la cartulina hacia abajo; la
\omterm{def:g5:air-pressure-weather:pressure}{atmósfera} la empuja hacia arriba. Gana la \omterm{def:g5:air-pressure-weather:pressure}{atmósfera} --- la
cartulina se queda y el agua se queda dentro.
\end{solution}

\begin{solution}{exo:g5:air-pressure-weather:5}
Lo que quitas es el \emph{\omterm{def:g2:air-around-us:air}{aire}} de dentro de la pajita. La
\omterm{def:g5:air-pressure-weather:pressure}{atmósfera}, que aprieta el zumo del vaso, empuja entonces el zumo por
la pajita despejada: quien hace el levantamiento es el cielo.
\end{solution}

\begin{solution}{exo:g5:air-pressure-weather:6}
Apretar la ventosa contra la pared expulsó el \omterm{def:g2:air-around-us:air}{aire} que había
detrás. Como ya no hay nada que empuje de vuelta desde debajo, el
empuje de la \omterm{def:g5:air-pressure-weather:pressure}{atmósfera} clava la ventosa a los azulejos.
\end{solution}

\begin{solution}{exo:g5:air-pressure-weather:7}
Se debilita --- queda menos océano de \omterm{def:g2:air-around-us:air}{aire} por encima que apriete
hacia abajo. Señales: los oídos que se destaponan en una subida
rápida y la falta de \omterm{def:g2:air-around-us:air}{aire} en un \omterm{def:g2:air-around-us:air}{aire} enrarecido (y el té que hierve
demasiado frío).
\end{solution}

\begin{solution}{exo:g5:air-pressure-weather:8}
El empuje de la \omterm{def:g5:air-pressure-weather:pressure}{atmósfera} --- la \omterm{def:g5:air-pressure-weather:pressure}{presión atmosférica}. Una bajada
rápida avisa de que viene una tormenta: el picnic, a cubierto.
\end{solution}

\begin{solution}{exo:g5:air-pressure-weather:9}
Cerrada en la cumbre, la botella guardó dentro \omterm{def:g2:air-around-us:air}{aire} enrarecido de
montaña; en el valle el empuje de fuera se hizo más fuerte que el
empuje encerrado dentro, y la botella se plegó. Llevada al revés, una
botella cerrada en el valle guarda dentro \omterm{def:g2:air-around-us:air}{aire} \emph{fuerte} mientras
el empuje de fuera se debilita con la altura: la botella se hincha
tensa --- y una bolsa de patatas fritas cerrada hace lo mismo en
cualquier viaje a la montaña.
\end{solution}

\begin{solution}{exo:g5:air-pressure-weather:10}
El \omterm{def:g2:air-around-us:air}{aire} fluye sin parar desde las pesadas regiones de presión alta
hacia las ligeras regiones de presión baja, y ese fluir es el
\omterm{def:g2:air-around-us:wind}{viento}. Así que la vieja definición se completa a sí misma: el
\omterm{def:g2:air-around-us:wind}{viento} es \omterm{def:g2:air-around-us:air}{aire} en movimiento --- movido por diferencias de
\omterm{def:g5:air-pressure-weather:pressure}{presión atmosférica}.
\end{solution}

\begin{solution}{exo:g5:air-pressure-weather:11}
Con la altura, el empuje de la \omterm{def:g5:air-pressure-weather:pressure}{atmósfera} sobre el agua se
debilita, así que hervir resulta más fácil: el \omterm{def:g4:changes-of-state:boiling-point}{punto de ebullición}
se desliza bastante por debajo de \qty{100}{\celsius}. Y, por la ley
de la meseta, el agua hirviendo no puede subir nunca por encima de su
\omterm{def:g4:changes-of-state:boiling-point}{punto de ebullición} --- la llama de más solo fabrica vapor. Así que
el agua de las cumbres hierve fría \emph{y se queda ahí}: té flojo,
pasta dura, física satisfecha.
\end{solution}
